%===============================================================================
% LLMWiki Technical Specification - Volume 10: Patent Draft
%===============================================================================
\chapter{Patent Draft}
\label{cha:patent}

This volume presents the patent claims for the novel aspects of LLMWiki. Claims are written in USPTO format with supporting specification references. Each claim corresponds to an architectural invariant or novel combination identified in Volumes 01-09.

\section{Field of the Invention}
\label{sec:patent:field}

The present invention relates to systems and methods for transforming heterogeneous digital artifacts into an evidence-aware semantic knowledge space, specifically a local-first semantic operating layer that enforces evidence closure through mandatory cross-encoder verification.

\section{Background}
\label{sec:patent:background}

Prior art includes Retrieval-Augmented Generation (RAG) systems (Lewis et al., 2020), GraphRAG (Edge et al., 2024), HippoRAG (Guti{\'e}rrez et al., 2024), LightRAG (Guo et al., 2024), and RAPTOR (Sarthi et al., 2024). These systems share common limitations: (1) vector indices as primary knowledge representation, (2) optional or post-hoc evidence verification, (3) dependence on external LLM APIs for critical-path operations, (4) lack of deterministic replay, and (5) no content-addressed unit identity.

\section{Summary of the Invention}
\label{sec:patent:summary}

The invention provides a semantic operating layer comprising: a content-addressed filesystem-backed write-ahead log; a property knowledge graph as canonical representation; a hybrid semantic index fusing vector, keyword, and graph signals; an evidence gate enforcing cross-encoder verification; and a DAG-based query planner with cost-based optimization.

\section{Independent Claims}
\label{sec:patent:independent}

\begin{claim}[System Architecture]
\label{claim:system}
A semantic operating layer system comprising:
\begin{enumerate}
  \item a filesystem-backed write-ahead log storing content-addressed semantic units, each unit identified by a SHA-256 hash of its canonical representation;
  \item a property knowledge graph wherein nodes represent semantic units and entities, and edges represent typed relationships, said graph persisted in a column-family store with prefix-encoded keys enabling type-partitioned scans;
  \item a hybrid semantic index combining a hierarchical navigable small world (HNSW) vector index, a BM25 inverted index, and a graph adjacency index, all keyed by the content-addressed unit identifiers;
  \item an evidence gate configured to receive a query and a candidate set of units from the hybrid index, and to verify each candidate by computing a cross-encoder relevance score between the query and the unit content, retaining only candidates whose score exceeds a calibrated threshold; and
  \item a query planner that decomposes an input query into a directed acyclic graph of subqueries, each subquery assigned a retrieval strategy, and executes said subqueries in topological order with cost-based optimization.
\end{enumerate}
\end{claim}

\begin{claim}[Method of Deterministic Ingestion and Replay]
\label{claim:replay}
A computer-implemented method for deterministic knowledge ingestion and replay, comprising:
\begin{enumerate}
  \item parsing a source artifact into a hierarchical tree of semantic units using a language-aware parser;
  \item computing a content address for each unit as SHA-256(canonical(content, type, span, metadata, children));
  \item extracting entity mentions and relations from units via rule-based patterns or a local language model;
  \item resolving entities against the knowledge graph using a weighted similarity function combining fuzzy name match, type compatibility, context overlap, and embedding cosine similarity;
  \item writing a diff of (added, modified, removed) unit identifiers to the filesystem-backed write-ahead log as an atomic append-only JSONL record;
  \item applying the diff to the knowledge graph and hybrid index within a single atomic transaction;
  \item wherein replay is achieved by sequentially applying logged diffs to an empty store, producing bitwise-identical state.
\end{enumerate}
\end{claim}

\par\noindent\emph{Prior art note (informal, see \S\ref{sec:patent:prior-art-notes}):} the content-identity/occurrence-lineage separation recited in Claims~\ref{claim:system}--\ref{claim:replay} (a SHA-256 content address decoupled from the location(s) at which that content occurs) is not independently novel; content-addressable storage systems with derivation/lineage tracking are directly anticipated by prior art (e.g.\ US8,135,760~B1, ``Determining the lineage of a content unit on an object addressable storage system,'' and related CAS patents). The identity split should be argued as one element of the broader system combination, not as a standalone novel mechanism.

\begin{claim}[Evidence Gate with Calibrated Threshold]
\label{claim:evidence-gate}
An evidence verification apparatus comprising:
\begin{enumerate}
  \item a cross-encoder model that takes a (query, unit) pair and outputs a relevance score in [0, 1];
  \item a calibrator trained on a validation set of (query, unit, label) triples, mapping raw scores to calibrated probabilities via isotonic regression or temperature scaling;
  \item a threshold selector that chooses an operating point $\tau$ such that false discovery rate $\leq \alpha$ on held-out data;
  \item a gate function $\mathcal{G}(q, \mathcal{R}, \tau) = \{ u \in \mathcal{R} \mid \text{calibrated\_score}(q, u) \geq \tau \}$;
  \item wherein the gate is a mandatory pipeline stage: no unit reaches the synthesis stage without passing the gate; and
  \item wherein the gate may return an empty set, triggering an "insufficient evidence" response rather than hallucinated synthesis.
\end{enumerate}
\end{claim}

\par\noindent\emph{Prior art note (informal, see \S\ref{sec:patent:prior-art-notes}):} the closest reference identified to date for a mandatory, threshold-gated cross-encoder stage ahead of generation is US20240070403A1 (``Grounded dialogue generation with cross-encoding re-ranker, grounding span prediction, and passage dropout''), which uses a cross-encoding re-ranker plus a grounding-span auxiliary task in a retrieve/re-rank/generate dialogue pipeline. It does not appear to gate to an empty set with an explicit ``insufficient evidence'' refusal (item 6 above); that refusal-over-hallucination behavior is the element most worth emphasizing when distinguishing this claim, but the comparison has not been confirmed against the reference's actual claim language and needs attorney review.

\begin{claim}[Deployment Architecture]
\label{claim:deployment}
A deployment method for a semantic operating layer, comprising:
\begin{enumerate}
  \item packaging the entire system---parsers, embedders, cross-encoder, graph store, indices, planner, runtime---as a single statically-linked binary with embedded ONNX models and tree-sitter grammars;
  \item storing all mutable state in a single directory \texttt{.llmwiki/} containing subdirectories for units, graph partitions, index files, and write-ahead log;
  \item enabling air-gapped operation by eliminating all external network dependencies in the critical path (ingestion, indexing, retrieval, verification);
  \item supporting horizontal scaling via directory synchronization (\texttt{rsync}, \texttt{git}, or object store replication) without consensus protocols;
  \item providing deterministic replay by copying \texttt{.llmwiki/} to a fresh machine and restarting the binary.
\end{enumerate}
\end{claim}

\begin{claim}[Label-Free Significance Calibration]
\label{claim:far-calibration}
A method for establishing a verification threshold in an evidence-gated retrieval system without relevance annotations, comprising:
\begin{enumerate}
  \item constructing a plurality of query--unit pairs for which genuine relevance is excluded by construction, said exclusion effected by at least one of: drawing the unit from a corpus declared disjoint in domain from that of the query; drawing the unit from a write-ahead-log snapshot predating creation of every entity referenced by the query; or drawing the unit from a graph component containing no entity mentioned in the query;
  \item evaluating the verification scorer over said pairs to accumulate an empirical null distribution of the scoring statistic;
  \item selecting as threshold the least score whose exceedance probability under said distribution, multiplied by the mean candidate count per query, attains a target false alarm rate; and
  \item emitting, with each admitted unit, said false alarm rate as the unit's significance, in place of a raw scorer output;
\end{enumerate}
wherein recalibration is triggered by advancement of the write-ahead log and requires no human annotation.
\end{claim}

\begin{claim}[Cascaded Proofreading Verification]
\label{claim:proofreading}
An evidence verification apparatus comprising a cascade of $n \geq 2$ verification stages, wherein:
\begin{enumerate}
  \item each stage evaluates a candidate against an evidence channel distinct from those of the other stages, said channels comprising at least lexical overlap, learned cross-encoder relevance, knowledge-graph reachability among co-surviving candidates, and inferential entailment;
  \item each stage irreversibly discards candidates scoring below a stage-specific threshold, discarded candidates being unrecoverable by subsequent stages;
  \item the stages are ordered by non-decreasing unit evaluation cost; and
  \item the false-acceptance rate of the cascade is the product of the per-stage false-acceptance rates, while the expected evaluation cost per candidate is the survival-weighted sum of per-stage costs.
\end{enumerate}
\end{claim}

\par\noindent\emph{Prior art note (informal, see \S\ref{sec:patent:prior-art-notes}):} the bare mechanism of a cascade of stages with increasing cost that irreversibly discards candidates is directly anticipated by the Viola--Jones cascade classifier (P.\ Viola and M.\ Jones, ``Rapid Object Detection using a Boosted Cascade of Simple Features,'' CVPR 2001; corresponding US patent filings on cascaded object detection). This claim should not rest on the cascade-with-irreversible-discard structure alone; it needs narrowing to the specific combination of RAG-evidence-verification channels recited in item 1 (lexical overlap, cross-encoder relevance, graph reachability, entailment) as applied to retrieval evidence candidates, which has no identified antecedent, rather than to cascading and early-discard as a general classification strategy.

\begin{claim}[Topological Consistency Gate]
\label{claim:consistency}
A method for certifying mutual consistency of a verified evidence set prior to answer synthesis, comprising:
\begin{enumerate}
  \item constructing a graph whose vertices are the verified units and whose edges join units sharing an entity or standing in a knowledge-graph relation;
  \item assigning to each vertex a stalk of admissible truth assignments over atomic claims extracted from the corresponding unit, to each edge a stalk over the claims shared across that edge, and to each incidence a restriction map induced by inferential entailment between said claim sets, thereby forming a cellular sheaf;
  \item computing the zeroth and first cohomology of said sheaf;
  \item admitting the evidence set for synthesis when a global section supporting the answer exists; and
  \item when the first cohomology is nonzero, withholding synthesis and emitting the support of the obstruction class as a machine-readable conflict witness identifying the relations carrying the inconsistency;
\end{enumerate}
whereby cyclic inconsistencies undetectable by pairwise contradiction testing are detected and localised.
\end{claim}

\par\noindent\emph{Prior art note (informal, see \S\ref{sec:patent:prior-art-notes}):} of the independent claims, this is the one for which searches to date found no directly competing prior art (existing RAG hallucination-detection work uses graph-, embedding-, or attention-based consistency signals, not sheaf cohomology). It is also the riskiest to draft: as a claim centered on a mathematical construction (cohomology of a cellular sheaf), it is more exposed to an abstract-idea objection (Alice/Mayo) than the other independent claims, and this exposure has not yet been analyzed. Recommend leading patent-attorney review with this claim.

\section{Dependent Claims}
\label{sec:patent:dependent}

\begin{claim}[Parser Extensibility]
\label{claim:parser-ext}
The system of Claim 1, further comprising a parser registry that dynamically loads parser implementations conforming to a \texttt{Parser} trait, each parser declaring its supported MIME types and output unit schema, whereby new languages and formats are added without modifying the core ingestion pipeline.
\end{claim}

\begin{claim}[Tree-sitter Integration]
\label{claim:tree-sitter}
The system of Claim 6, wherein at least one parser uses a tree-sitter grammar to produce a concrete syntax tree, and semantic units are extracted via S-expression query patterns (SCM) mapped to unit fields, enabling incremental re-parsing with $O(\log n)$ edit complexity.
\end{claim}

\begin{claim}[Hierarchical Chunking]
\label{claim:chunking}
The system of Claim 1, further comprising a chunker that recursively partitions semantic units along structural boundaries (functions, classes, sections) while respecting a maximum token budget and configurable overlap, producing chunks that preserve hierarchical context via parent unit references.
\end{claim}

\begin{claim}[Local Embedding Inference]
\label{claim:local-embed}
The system of Claim 1, wherein the embedder executes entirely on local hardware (CPU, GPU, or NPU) using ONNX Runtime with model quantization (INT8, INT4), supporting batch inference with dynamic padding, and requiring zero external API calls.
\end{claim}

\begin{claim}[Entity Resolution]
\label{claim:entity-resolution}
The system of Claim 1, wherein entity resolution merges mentions into canonical entities using a scoring function $S(e_1, e_2) = \alpha \cdot \text{Fuzzy}(n_1, n_2) + \beta \cdot \mathbb{1}[t_1 = t_2] + \gamma \cdot \text{Jaccard}(ctx_1, ctx_2) + \delta \cdot \cos(emb_1, emb_2)$, with $\alpha+\beta+\gamma+\delta=1$, and creates a new entity iff $\max S < \tau_{er}$.
\end{claim}

\begin{claim}[Graph Adjacency Index]
\label{claim:adj-index}
The system of Claim 1, wherein the graph adjacency index is a separate column family storing for each unit a serialized list of (relation\_type, target\_unit\_id, weight) tuples, enabling sub-millisecond $k$-hop expansion during hybrid retrieval without accessing the primary graph store.
\end{claim}

\begin{claim}[Hybrid Retrieval with Reciprocal Rank Fusion]
\label{claim:hybrid-rrf}
The system of Claim 1, wherein hybrid retrieval combines vector, keyword, and graph scores via reciprocal rank fusion: $score(u) = \sum_{s \in \{v,k,g\}} w_s \cdot \frac{1}{k + rank_s(u)}$, where $w_s$ are configurable weights and $rank_s(u)$ is the position of $u$ in the top-$K$ results of signal $s$.
\end{claim}

\begin{claim}[DAG Query Planner]
\label{claim:dag-planner}
The system of Claim 1, wherein the query planner produces a plan represented as a directed acyclic graph $P = (V, E)$ where each vertex $v \in V$ is a retrieval, synthesis, or tool invocation step with declared inputs and outputs, and edges $E$ represent data dependencies, and the executor schedules steps in topological order with parallel fan-out where dependencies permit.
\end{claim}

\begin{claim}[Budgeted Execution]
\label{claim:budget}
The system of Claim 10, further comprising a budget tracker that enforces hard limits on cumulative tokens, wall-clock time, tool invocations, and monetary cost, with configurable overflow behavior: truncate, return partial, or error.
\end{claim}

\begin{claim}[Streaming Protocol]
\label{claim:streaming}
The system of Claim 1, further comprising a Server-Sent Events (SSE) endpoint that emits fine-grained events: plan\_start, step\_start, evidence\_found, token, step\_complete, plan\_complete, budget\_warning, enabling real-time UI rendering of the query execution pipeline.
\end{claim}

\begin{claim}[Evaluation Framework]
\label{claim:eval}
A computer-implemented evaluation method for a semantic operating layer, comprising:
\begin{enumerate}
  \item datasets with human-annotated evidence spans for each query;
  \item metrics: Evidence Recall@k, Evidence Precision@k, Evidence F1@k, Gate Rejection Rate, Attribution Score (claim-level NLI verification against citations), and Citation Quality (mean cross-encoder score of cited units);
  \item automated configuration sweeps over retrieval $k$, hybrid weights, gate threshold $\tau$, embedder model, and LLM model;
  \item statistical rigor: bootstrap confidence intervals, paired Wilcoxon tests, Benjamini-Hochberg correction, effect size reporting;
  \item regression gates in continuous integration blocking merges on $>2\%$ metric degradation.
\end{enumerate}
\end{claim}

\section{Design Decisions Supporting Novelty}
\label{sec:patent:decisions}

\begin{table}[htbp]
\centering
\caption{Design Decisions Supporting Novelty}
\label{tab:patent:decisions}
\begin{adjustbox}{max width=\textwidth}
\begin{tabular}{lll}
\toprule
\textbf{Decision} & \textbf{Novelty Argument} & \textbf{Prior Art Gap} \\
\midrule
Filesystem-backed event log & No RAG system uses directory sync for replay & All require Kafka/Redis/DB WAL; general append-only-log durability is well established outside RAG (\S\ref{sec:patent:prior-art-notes}) \\
Graph as canonical query representation & Vectors/keyword index are projections, not primary & RAG: vectors primary; GraphRAG: LLM-built graph, not rebuildable from a log \\
Mandatory evidence gate & Not a reranker; zero results allowed & Closest identified reference (US20240070403A1) reranks and grounds but does not appear to gate to empty-set refusal; not fully confirmed \\
Content identity / occurrence lineage split & SHA-256 content identity decoupled from location identity & Not independently novel: anticipated by CAS lineage-tracking prior art (e.g.\ US8,135,760~B1); supports Claim~\ref{claim:system} only as part of the full combination \\
Local-first, LLM-optional & Critical path has no API calls & All major RAG frameworks require LLM API on critical path \\
DAG planner with costs & Topological execution with parallelism & RAG: single retrieval; GraphRAG: fixed patterns \\
Deterministic replay & Bitwise identical re-execution & No directly competing RAG system found; general event-sourcing replay is well established (not RAG-specific prior art found) \\
Cohomological consistency gate & Detects cyclic inconsistency no pairwise checker can localise & No direct competing prior art found to date (\S\ref{sec:patent:prior-art-notes}); not checked against patent databases specifically; separately exposed to abstract-idea (Alice/Mayo) objection \\
\bottomrule
\end{tabular}
\end{adjustbox}
\end{table}

\section{Specification Support}
\label{sec:patent:spec-support}

Each claim is supported by the specification in Volumes 01-09:
\begin{itemize}
  \item Claims 1-5: Vol. 01 (\S\ref{sec:architecture:three-layer}, \ref{sec:architecture:filesystem-wal}, \ref{sec:architecture:invariants}), Vol. 02 (\S\ref{sec:parsing:incremental}), Vol. 03 (\S\ref{sec:kg:entity-resolution}, \ref{sec:kg:index}), Vol. 04 (\S\ref{sec:retrieval:evidence-gate}), Vol. 07 (\S\ref{sec:runtime:plan-structure})
  \item Claims 6-8: Vol. 02 (\S\ref{sec:parsing:parser-trait}, \ref{sec:parsing:tree-sitter}, \ref{sec:parsing:chunking}), Vol. 02 (\S\ref{sec:parsing:embedding})
  \item Claims 9-10: Vol. 03 (\S\ref{sec:kg:entity-resolution}, \ref{sec:kg:adjacency-index}, \ref{sec:kg:hybrid-search})
  \item Claims 11-13: Vol. 07 (\S\ref{sec:runtime:plan-structure}, \ref{sec:runtime:budget}, \ref{sec:runtime:streaming})
  \item Claim 14: Vol. 06 (\S\ref{sec:eval:metrics}, \ref{sec:eval:pipeline}, \ref{sec:eval:statistical})
\end{itemize}

\section{Prior Art Research Notes (Informal)}
\label{sec:patent:prior-art-notes}

\begin{quote}
\textbf{Status: informal research, not legal advice.} The notes in this section come from web and academic-literature search (Google Patents free-text search, arXiv, general web search), conducted as a sanity check before any filing decision, not a professional patentability or freedom-to-operate opinion. No paid patent-database search (e.g.\ full-text USPTO/EPO with a patent search firm) has been run. \textbf{Do not file, disclose publicly, or rely on the ``non-obvious'' language elsewhere in this volume without review by qualified patent counsel.} If public disclosure occurs before a filing decision is made, note that South Korea provides a 12-month grace period for inventor-originated public disclosure (this exception's exact procedural requirements have not been independently verified here and must be confirmed with counsel before relying on it).
\end{quote}

\begin{itemize}
  \item \textbf{Cascaded, cost-ordered, irreversible-discard verification} (Claim~\ref{claim:proofreading}): closely anticipated as a general mechanism by the Viola--Jones cascade classifier (CVPR 2001) and its associated patent filings, which cascade increasingly expensive stages and irreversibly reject candidates. The RAG-specific channel combination (lexical, cross-encoder, graph-reachability, entailment) is the defensible remainder; the cascade-with-discard structure itself is not new.
  \item \textbf{Content-addressable identity with lineage/occurrence tracking} (Claims~\ref{claim:system}, \ref{claim:replay}): anticipated by object-addressable/content-addressable storage patents tracking content-unit lineage and derivation, notably US8,135,760~B1 (``Determining the lineage of a content unit on an object addressable storage system''), and by adjacent CAS patents (e.g.\ US8,335,890~B1 ``Associating an identifier with a content unit''; US8,909,875~B1; US8,560,503~B1). Content hashing plus a separate location/lineage record is an established pattern (also visible in Git's own object model), not a novel mechanism in isolation.
  \item \textbf{Mandatory, threshold-gated cross-encoder verification ahead of generation} (Claims~\ref{claim:system}, \ref{claim:evidence-gate}): the closest reference found is US20240070403A1 (``Grounded dialogue generation with cross-encoding re-ranker, grounding span prediction, and passage dropout''), a retrieve/re-rank/generate dialogue pipeline using a cross-encoding re-ranker and a grounding-span auxiliary task. Also adjacent: US20240070403A1's sibling RAG-optimization filings (e.g.\ patents on tuning cross-encoder/bi-encoder/LLM-ranker combinations in retriever pipelines) and multiple 2024--2026 RAG patent filings using cross-encoder reranking. None found so far recite gating to an explicit empty-set ``insufficient evidence'' refusal as opposed to always returning top-$K$ reranked results, but this distinction has not been checked against the actual claim language of those filings.
  \item \textbf{Sheaf-cohomology-based evidence consistency checking} (Claim~\ref{claim:consistency}): no directly competing prior art identified in this round of searching. Existing RAG hallucination/inconsistency detection approaches found in literature search (evidence-graph consistency methods, semantic-entropy methods, attention/embedding-based contradiction detectors) use different mathematical machinery. This is the strongest remaining novelty candidate based on searches to date, but (a) no dedicated patent-database search has been run for it specifically, and (b) a claim built around a cohomological computation is more exposed to an abstract-idea (Alice/Mayo, 35 U.S.C. \S 101) rejection than the other independent claims, which this research did not evaluate.
  \item \textbf{Label-free / empirical-null threshold calibration} (Claim~\ref{claim:far-calibration}): related academic technique (score-distributional threshold optimization without relevance judgments, cf.\ Arampatzis et al.) exists, and empirical null-hypothesis testing is well-established general statistics; no directly competing patent found for the specific corpus/temporal/graph-displacement construction used here to build the null. Likely a defensible narrow claim if limited to that specific construction rather than empirical-null calibration in general.
  \item \textbf{Approval-gated document mutation with automatic knowledge-graph reindexing} (Volume 12, not yet claimed in this volume): adjacent prior art exists for document approval workflows generally (e.g.\ US20230316190 ``Document workflows in a document management system''; US10,630,858 ``Document approval management system''), a crowded space. No reference found combining an approval gate with atomic re-projection into a semantic knowledge graph specifically, but this was a narrower, lower-priority search and should be re-run before any claim is drafted for Volume 12's mechanism.
\end{itemize}

\section{Potential Design-Arounds and Defense}
\label{sec:patent:defense}

\begin{itemize}
  \item \textbf{Design-around}: Use vector DB with post-hoc reranking. \textbf{Defense}: Claim 3 requires mandatory gate before synthesis; reranking is not gating.
  \item \textbf{Design-around}: Use database WAL instead of filesystem. \textbf{Defense}: Claim 1 requires filesystem-backed WAL with directory-sync replay.
  \item \textbf{Design-around}: Skip graph, use only hybrid vector/keyword. \textbf{Defense}: Claim 1 requires graph as canonical knowledge representation with adjacency index.
  \item \textbf{Design-around}: Cloud-only deployment. \textbf{Defense}: Claim 4 requires single-binary air-gapped operation.
  \item \textbf{Design-around}: Calibrate the threshold on a labelled relevance set, avoiding the null construction. \textbf{Defense}: Claim~\ref{claim:far-calibration} is directed to the label-free case specifically; a competitor taking this route incurs the annotation cost the claim exists to eliminate, and cannot recalibrate automatically on log advancement.
  \item \textbf{Design-around}: Use multiple rerankers in sequence without irreversible discard, retaining all candidates and combining scores. \textbf{Defense}: Claim~\ref{claim:proofreading} recites irreversible discard, channel distinctness, and cost ordering jointly; a score-combining ensemble obtains neither the multiplicative error suppression nor the survival-weighted cost bound, and is a different mechanism rather than a variant.
  \item \textbf{Design-around}: Detect contradictions pairwise between retrieved passages. \textbf{Defense}: Claim~\ref{claim:consistency} is directed to cohomological detection, which is strictly stronger --- cyclic inconsistency admits no pairwise witness, so a pairwise checker does not read on the claim and does not achieve its effect.
\end{itemize}

\paragraph{On the ``combination of known elements'' objection.} The independent claims of \S\ref{sec:patent:independent} are vulnerable to the argument that hybrid retrieval, cross-encoders, and graph stores are individually well known. Claims~\ref{claim:far-calibration}--\ref{claim:consistency} were drafted to sit outside that objection by reciting a construction procedure with a stated technical effect rather than an arrangement of components; that intent still stands, but the informal search in \S\ref{sec:patent:prior-art-notes} found that two of the claimed elements do have antecedents outside retrieval specifically: the cascade-with-irreversible-discard structure of Claim~\ref{claim:proofreading} (Viola--Jones) and the content-identity/lineage split underlying Claim~\ref{claim:system} (CAS lineage patents). The displacement conditions of Claim~\ref{claim:far-calibration} and, most notably, the sheaf construction of Claim~\ref{claim:consistency} remain without an identified antecedent as of this search round and are the strongest candidates for the non-obviousness argument; Claims~\ref{claim:proofreading} and~\ref{claim:system} should be narrowed to their RAG-specific elements (the channel combination; the combination with graph/index rebuild from the log) rather than argued on the cascade or content-addressing mechanisms alone. This assessment is informal and requires confirmation by patent counsel before it is relied on.

%===============================================================================
% End of Volume 10
%===============================================================================